The entire esophageal wall is part of the gastrointestinal (GI) tract. It follows the standard four-layer GI architectural plan (mucosa, submucosa, muscularis propria, adventitia/serosa). Its inner lining is non-keratinized stratified squamous epithelium (adapted for abrasion from swallowed food), which transitions abruptly at the Z-line (esophagogastric junction) into the simple columnar epithelium of the stomach. 

Fundic distension itself (from gas, a large meal, delayed emptying, etc.) is a physiological trigger, not a disease. Healthy people experience gastric distension and occasional TLESRs every day. The distinction between normal and pathological lies in:the frequency, duration, or inappropriate timing of barrier relaxations,
the volume or composition of material that refluxes,
the inability of the esophagus to clear it,
and the resulting mucosal injury, hypersensitivity, or symptoms.

Distension is the common initiating mechanical stimulus in the pathway; it is not the disease entity.

Summary of the distinction
Neural (vagus + phrenic) coordination failure can be an important mechanism.
The esophagus is GI tissue; the barrier is the esophagogastric junction; the injurious material comes from the stomach.
Therefore the resulting clinical condition is classified and treated as a gastroesophageal disorder, not as a primary stomach-acid disease or as a primary neurological disease.

The hypothesis correctly highlights under-appreciated neuromuscular aspects. 
